\section*{対象読者と読む順序}
二次関数，整式の計算，三角関数の基本，指数・対数，数列の和を学んだ読者を想定する。
微積分を初めて学ぶ場合は，\chapterref{章:導入}から\chapterref{章:積分の応用}までを順に読み，章末問題で計算を確かめてほしい。
\chapterref{章:極限と微分}は後の説明に必要な微分をまとめているため，他の章より多くの内容を含む。
すでに微分を学んでいる場合は，\chapterref{章:定積分}から読み始め，必要になった箇所だけ\chapterref{章:極限と微分}へ戻ってよい。

\begin{center}
\begin{tabularx}{\linewidth}{@{}p{35mm}X@{}}
\toprule
読みたい内容 & 主な箇所\\
\midrule
積分の意味を知る & 第1〜3章。近似の和が定積分になり，原始関数で計算できる理由を扱う。\\
計算方法を選ぶ & 第4〜7章。式の整理，置換，部分積分，対称性を比較する。\\
積分を問題に使う & 第8〜9章。図形・運動・数列・不等式を通して，何を足すかを考える。\\
手法の背景へ進む & 第10〜11章。置換の幾何，逆関数，近似，広義積分を扱う。\\
自分で方針を決める & \chapterref{章:総合演習}。問題を解いてからヒントや巻末の解答を参照する。\\
\bottomrule
\end{tabularx}
\end{center}

\chapterref{章:積分と数列}の一般のベータ関数，第10〜11章，および\chapterref{章:積分の応用}の重心とパップス＝ギュルダンの定理は，高校の基本事項から先へ進むための発展内容である。
基本の学習に際して，これらの定理をすべて覚える必要はない。
発展内容では，高校の計算から何が広がるかと，新たに何を仮定するかを区別して読むとよい。
本書は高校数学全分野の教科書ではない。複素数平面，確率，整数などは，積分との関係が必要な範囲で扱う。

